\documentclass[10pt,a4paper]{amsart}
\usepackage[margin=19mm]{geometry}
\usepackage{amsmath,amssymb,mathtools}
\usepackage[scr=rsfs,cal=euler]{mathalfa}
\usepackage{xfrac}
\usepackage[unicode,hidelinks,bookmarks=false]{hyperref}
\newcommand{\ie}{i.e.\ }
\newcommand{\cf}{cf.\ }
\newcommand{\resp}{respectively\ }
\newcommand{\Subsection}{\subsection}
\providecommand{\nobreakdash}{\nobreak\hbox{-}}
\setlength{\emergencystretch}{2em}
\multlinegap0pt
\usepackage{xcolor}
\usepackage{amssymb}
\usepackage{enumerate}
\usepackage[shortlabels]{enumitem}
\usepackage{mathtools}
\newtheorem{corollary}{Corollary}
\newtheorem{proposition}{Proposition}
\newcommand{\Id}{\mathrm{id}}
\title{Measure preserving maps with bounded total variation}
\author{Stefano Bianchini, Luca Talamini}
\date{Source: arXiv:2603.18819v1 (CC BY 4.0)}
\begin{document}
\maketitle

\noindent\emph{Source and licence.} Measure preserving maps with bounded total variation. Version arXiv:2603.18819v1 (CC BY 4.0), distributed by arXiv under the Creative Commons Attribution 4.0 International licence, http://creativecommons.org/licenses/by/4.0/, as recorded in the arXiv licence metadata for this exact version and shown on the abstract page https://arxiv.org/abs/2603.18819v1. Full licence notice: \texttt{licenses/arxiv-2603.18819-CC-BY-4.0.txt}.\par\smallskip

\noindent\textit{Editorial scope.} Counted unit: the article's proposition Prop:laplacian\_1 (source0.tex lines 750--756). The article's complete original proof of it is delivered with it (source0.tex lines 758--790, one \texttt{proof} environment of 33 lines). No mathematics is written, repaired or re-derived: the statement and the proof are the article's own lines, in the article's own notation, and no retained line is substituted or re-flowed.  Retained uncounted as the source-local context the proof needs: lines 736--746.  Nothing was omitted from any retained span, so the package carries no omission record.  No retained line contains a citation, so the wrapper carries no bibliography and no bibliography entry.  No retained line contains a figure environment, an image-inclusion command or drawing code, so no image is delivered and the package declares no asset dependency.  Editorial formatting only, in addition: an AMS \texttt{amsart} preamble that restates the article's own declarations and macros used by the retained lines, namely the environments \texttt{corollary}, \texttt{proposition} on separate counters, together with the standard packages those lines need.  The retained environments print in this document as Corollary 1, Proposition 1 in order of appearance, whereas the article prints the same items with the numbering of its own full document.  The complete native source of the article as distributed in the arXiv e-print for this exact version is bundled unchanged as \texttt{sources/196726/source0.tex}.
\begin{corollary}
\label{Cor:conser_transprto}
Let $\phi \in \mathbf W^{1,2}(\Omega)$ such that \emph{(\textbf{MPC})} holds. Let $\nabla f_t, \Phi_t$ be the maps in the polar decomposition of $T_t$, i.e.
$$
T_t = \nabla f_t \circ \Phi_t.
$$
Then 
\begin{equation*}
\lim_{t \to 0} \bigg\| \frac{{\nabla f_t} - \Id}{t} - \nabla \phi \bigg\|_{\mathbf L^2} = 0, \quad \lim_{t \to 0} \bigg\| \frac{\Phi_t - \Id}{t} \bigg\|_{\mathbf L^2} = 0,
\end{equation*}
\end{corollary}

\begin{proposition}
\label{Prop:laplacian_1}
Let $\phi \in \mathbf W^{1,2}(\Omega)$ such that \emph{(\textbf{MPC})} holds. Then
$$
\Delta \phi \geq 0.
$$
\end{proposition}

\begin{proof}
Let $\nabla f_t, \Phi_t$ be as above.
Being $\nabla {f_t}$ measure preserving and monotone, hence BV, it holds
\begin{equation*}
\det \big( \Id + t D^a \nabla h_t(x) \big) = 1,
\end{equation*}
for $\mathscr L^d$-a.e. $x \in \Omega$, where we recall 
$$
\nabla h_t = \frac{\nabla f_t - \Id}{t}.
$$
If $\lambda_{i}(t,x)$ are the eigenvalues of $D^a \nabla h_t(x)$ one gets, using the arithmetic-geometric inequality
\begin{equation}
\label{Equa:eigenvalues_nabla_h_t_1}
1 = \bigg( \prod_i (1 + t \lambda_{i}(t,x)) \bigg)^{\frac{1}{d}} \leq \frac{1}{d} \sum_i (1 + t \lambda_{i}(t,x)) = 1 + \frac{t}{d} \sum_i \lambda_{i}(t,x).
\end{equation}
By \eqref{Equa:eigenvalues_nabla_h_t_1}, there holds
$$
\mathrm{Tr}(D^a \nabla h_t) \geq 0
$$
while for the singular part of the Laplacian, one has
\begin{equation*}
t (\Delta h_t)^s = (\Delta {f_t})^s \geq 0,
\end{equation*}
being $f_t$ convex. Hence $\Delta h_t \geq 0$ in the sense of distributions, 
$$
\int_{\Omega} \nabla \psi(x) \cdot \nabla h_t(x) \mathrm{d} x \leq 0 \qquad \text{for all $\psi \in C^1_c(\Omega)$, $\psi \geq 0$}.
$$
Since by Corollary \ref{Cor:conser_transprto} we have the convergence $\nabla h_t \to \nabla \phi$  in $\mathbf L^2$, we deduce for every $\psi \in C^1_c(\Omega)$ that 
$$
0 \geq \lim_{t \to 0^+} \int_{\Omega} \nabla \psi(x) \cdot \nabla h_t(x) \mathrm{d} x = \lim_{t \to 0^+} \int_{\Omega} \nabla \psi(x) \cdot \nabla \phi(x) \mathrm{d} x 
$$
and this proves the result.
\end{proof}

\end{document}
